\documentclass{report}
\usepackage[T1]{fontenc} 
\usepackage[utf8]{inputenc} 
\usepackage{lmodern}
\usepackage [frenchb]{babel}
% Pour pouvoir utiliser 
\usepackage{ucs}

\usepackage{textcomp}
\usepackage{graphicx}
\usepackage{keystroke}
\usepackage{amssymb}
\usepackage{amsmath}
%\renewcommand{\thesection}{\arabic{section}} % numérotation des sections	
\usepackage[cc]{titlepic} %rajouter le logo dans la page de garde
\usepackage{url} % Pour avoir de belles url
\usepackage[a4paper]{geometry}
\usepackage[linktocpage]{hyperref}
% Pour pouvoir faire plusieurs colonnes
\usepackage {multicol}
%Pour faire plusieurs lignes
\usepackage{multirow}
\usepackage{slashbox}
% Pour mettre du code source
\usepackage {listings}
% Pour pouvoir passer en paysage
\usepackage{lscape}	

% Pour pouvoir faire plusieurs colonnes
\usepackage {multicol}

% POur crééer un index
\usepackage{makeidx}

%insertion de pdf
\usepackage{pdfpages}


%Pour taper des algorithme en pseudo code
\usepackage{algorithm,algorithmic}
%Utilisation d'acronyme
\usepackage{acronym}

\usepackage{fancyhdr}
\pagestyle{fancy}
\usepackage{lastpage} %numérotation type k/n
\renewcommand{\footrulewidth}{0.5pt} % trait horizontal
%\renewcommand{\headrulewidth}{1pt} % suppresion du trait horizontal dans l'entête
\lfoot{Projet Services}
\cfoot{\thepage \hspace{0.15cm}sur  \pageref{LastPage} } % numérotation des pages
\rfoot{A.\textsc{M} \& R.\textsc{R}}
\fancyhead[L]{\hspace{-1cm}\includegraphics[scale=0.105]{img/logouniv.png}} %logo entête !


\fancypagestyle{plain}{% 1ères pages des chapitres:
  % \fancyhead{} % supprime l’entete...
  % \renewcommand{\headrulewidth}{0pt} % ...et le filet
}




\hypersetup{
  backref=true,
  %permet d'ajouter des liens dans...
  pagebackref=true,%...les bibliographies
  hyperindex=true, %ajoute des liens dans les index.
  colorlinks=true, %colorise les liens
  breaklinks=true, %permet le retour à la ligne dans les liens trop longs
  urlcolor= blue, %couleur des hyperliens
  citecolor=	cyan,
  bookmarks=true, %créé des signets pour Acrobat
  bookmarksopen=true,
  %si les signets Acrobat sont créés,
  %les afficher complètement.
  pdftitle={Rapport Projet Services}, %informations apparaissant dans
  pdfauthor={MARGUERITE Alain},
  %dans les informations du document
  pdfsubject={Doc}
  %sous Acrobat.
}

\makeindex
%%%% debut macro pour enlever le nom chapitre %%%%
\makeatletter
\def\@makechapterhead#1{%
  \vspace*{50\p@}%
          {\parindent \z@ \raggedright \normalfont
            \interlinepenalty\@M
            \ifnum \c@secnumdepth >\m@ne
            \Huge\bfseries \thechapter\quad
            \fi
            \Huge \bfseries #1\par\nobreak
            \vskip 40\p@
}}

\def\@makeschapterhead#1{%
  \vspace*{50\p@}%
          {\parindent \z@ \raggedright
            \normalfont
            \interlinepenalty\@M
            \Huge \bfseries  #1\par\nobreak
            \vskip 40\p@
}}
\makeatother
%%%% fin macro %%%%
%Couverture 
\widowpenalty=10000
\clubpenalty=10000
\hyphenation{permet-tant}

\title{ {\huge \textsc{Distributed System}} \\Master ALMA 2\up{eme} année \\\vspace{3cm} \emph{Causal Order Message Delivery : From Logical to Physical Time}\\{\small Encadrant : M.\textsc{Achour}}\vspace{3cm}}
\titlepic{ \includegraphics[scale=0.2]{img/logouniv}}
\author{A.\textsc{Marguerite}\\ R.\textsc{Rincé}\vspace{3cm}}
\date{Université de Nantes \\ 2 rue de la Houssinière, BP92208, F-44322 Nantes cedex 03, FRANCE}
%\vspace{20cm}


\newcommand{\cf}{\emph{cf}.}


\begin{document}

\maketitle
\renewcommand{\labelitemi}{$\bullet$} 

\clearpage

\tableofcontents
\clearpage

%\chapter*{Remerciements}
\phantomsection
\addcontentsline{toc}{section}{Remerciements}
 TODO TODO TODO TODO TODO TODO TODO TODO TODO TODO TODO TODO TODO TODO TODO TODO TODO TODO TODO TODO TODO TODO

 TODO TODO TODO TODO TODO TODO TODO TODO TODO TODO TODO TODO TODO TODO TODO TODO TODO TODO TODO TODO

%\chapter{Introduction}\label{chap:Intro}

\section{Problématiques}
Les problèmes d'inter-opérabilité et de réutilisabilité font partie des problématiques majeures des systèmes informatiques. La programmation objet ne parvient pas à répondre à ces problématiques à cause notamment de ses propriétés de granularité et de fort couplage. Les problèmes de cohésions (pas d'explications des besoins requis) sont un frein à l'utilisation de la programmation objet à grande échelle. 

Parmi les solutions à ces problématiques les architectures à composants sont des alternatives efficaces. Un composant, entité principale, correspond par définition à une ou plusieurs fonctionnalités. Il est donc possible de le concevoir et de le pré-tester de manière indépendante. La coopération entres composant est alors plus aisée. De plus un composant possède un couplage lâche. Cette propriété garantissant l'inter-opératibilté est permise par l'utilisation d'interfaces entre composants. 

\section{Présentation du sujet}
Le module \textsc{Architectures and Components} nous propose une étude de la programmation par composants. L'objectif de ce Projet/TP  est de réaliser un Home Architecture Description Language, et d'implémenter une architecture classique \emph{Clients-Serveur}. La démarche proposée pour atteindre cet objectif est composée de trois étape : 


\begin{enumerate}
\item 
Réalisation du M2 :\hfill \\
Il s'agit d'un méta-modèle de l'architecture en composants.  
\item
Réalisation du M1 :\hfill \\
À partir du méta-modèle précédement réalisé, nous pouvons proposer une instance de ce dernier et obtenir un modèle. Nous choisirons ici de modéliser une application classique \emph{Client-Serveurs}.  
\item
Réalisation du M0 :\hfill \\
Dernière étape, elle consiste à l'implémentation du M1 et de coder concrétement qui y sont modélisées.
\end{enumerate}

Dans ce rapport nous détaillerons chaque étape la modélisation proposée puis l'implémentation.

\chapter{Logical Time}
\section{Question 1 Design a logical-time base algorithm}

Posons 
\begin{algorithm}
\caption{\textbf{CO\_Reception}( Message m\{Emetteur e, ListPred pred\})}
/* \textit{ où pred est la liste des prédécesseurs */}
\label{algo:recherche}
\begin{algorithmic}[1]
  \IF{ $\forall$ p$\in$ pred $\mid$ p.sn $\leq$ del(p.i)}
     \STATE  del(e.i) ++
     \STATE  cb +=\{e.sn, e.i\}
   \STATE  CO\_Deliver(m)
   \FORALL{ m\'~ $\in$ ListMessages}
   CO\_Receive(m\'~)
   \ENDFOR
   \ELSE
   \STATE ListMessages.add(m)
  \ENDIF
 \end{algorithmic}
\end{algorithm}


\section{Question 2 Analyse of the algorithm}

\begin{algorithm}
\caption{\textbf{CO\_Broadcast}()}
/* \textit{ Avec horloge globale */}
\label{algo:broad}
\begin{algorithmic}[1]
\STATE Broadcast(m,cb)
\STATE S_n $\leftarrow$ GC.clock
\STATE cb $\leftarrow $ \{(sn, i\}
\end{algorithmic}
\end{algorithm}



\begin{algorithm}
\caption{\textbf{CO\_Receive}( Message m\{Emetteur,ListPred\})}
/* \textit{ Avec horloge globale */}
\label{algo:q3}
\begin{algorithmic}[1]
  \IF{ $\forall$ p$\in$ cb $\mid$ p.sn $\leq$ del(p.i)}
   \STATE  del(e.i) $\leftarrow$ e.sn
   \STATE  CO\_Deliver(m)
   \FORALL{ m\'~ $\in$ ListMessages}
   CO\_Receive(m\'~)
   \ENDFOR
   \ELSE
   \STATE ListMessages.add(m)
  \ENDIF  
\end{algorithmic}
\end{algorithm}




\chapter{From Logical Time to Physical Time}

\section{Question 3 Modify your algorithm from logical time to physical time}

Posons les variables globales du système

\begin{itemize}
\item  \verb+GlobalClock+ : Horloge globale du système.
\end{itemize}


Redéfinissions la procédure \verb+CO_Brodcast+ : \ref{algo:broad}.

\begin{algorithm}
\caption{\textbf{COBroadcast}()}
/* \textit{ Avec horloge globale */}
\label{algo:broad}
\begin{algorithmic}[1]
\STATE Broadcast(m,cb)
\STATE S\_n $\leftarrow$ GC.clock
\STATE cb $\leftarrow $ \{(sn, i\}
\end{algorithmic}
\end{algorithm}



\begin{algorithm}
\caption{\textbf{CO\_Receive}(Message m)}
/* \textit{ Avec horloge globale */}
\label{algo:q3}
\begin{algorithmic}[1]
  \IF{ $\forall$ p$\in$ m.cb $\mid$ p.sn $\leq$ del(p.i)}
   \STATE  del(m.i) $\leftarrow$ m.sn
   \STATE  CO\_Deliver(m)
   \FORALL{ m\'~ $\in$ ListMessages}
   \STATE   CO\_Receive2(m')
   \ENDFOR
   \ELSE
   \STATE ListMessages.add(m)
  \ENDIF  
\end{algorithmic}
\end{algorithm}


\begin{algorithm}
\caption{\textbf{CO\_Receive2}(Message m)}
/* \textit{ Avec horloge globale */}
\label{algo:receive2}
\begin{algorithmic}[1]
  \IF{ $\forall$ p$\in$ m.cb $\mid$ p.sn $\leq$ del(p.i)}
   \STATE  del(m.i) $\leftarrow$ m.sn
   \STATE  CO\_Deliver(m)
   \STATE ListMessages.delete(m)
   \FORALL{ m' $\in$ ListMessages}
   \STATE CO\_Receive2(m)
   \ENDFOR
  \ENDIF  
\end{algorithmic}
\end{algorithm}

% LocalWords:  From Logical Time to Physical Modify your algorithm CO
% LocalWords:  from logical time physical GlobalClock Brodcast cb sn
% LocalWords:  Broadcast GC.clock Receive m.cb p.sn del p.i m.i m.sn
% LocalWords:  Deliver ListMessages aa ListMessages.add aaa
% LocalWords:  ListMessages.delete

\chapter{Messages with Bounded Lifetime}

\section{System model : messages with a bounded lifetime}

\subsection{Question 4 Extended the previous algorithm}

\begin{algorithm}
\caption{\textbf{Send\_Time}()}
/* \textit{ Avec horloge globale */}
\label{algo:Q4a}
\begin{algorithmic}[1]
  \IF{ $\exists$ M $\in$ ListeMessages  $\mid$ \{ m.send\_time + $\Delta$ = GC.clock() \}}
      \STATE  CO\_Deliver(M)
   \FORALL{ p $\in$ m.pred }
   \IF{del(p.i) $\leq$ p.sm}
   del(p.i) $\leftarrow$ p.sn
   \ENDIF
   \ENDFOR
  \ENDIF
\end{algorithmic}
\end{algorithm}


\begin{algorithm}
\caption{\textbf{CO\_Recceive\_With\_Delay}(Message m)}
/* \textit{ Au moment où on reçoit un message */}
\label{algo:Q4b}
\begin{algorithmic}[1]
  \IF{m.send\_time + $\Delta$ $\leq$ GC.Clock}
  \STATE Delete(m)
  \ELSE[$\forall$ p$\in$ pred $\mid$ p.sn $\leq$ del(p.i)]
   \STATE  del(e.i) $\leftarrow$ e.sn
     \STATE  cb +=\{e.sn, e.i\}
     \STATE  CO\_Deliver(m)
     \FORALL{m\'~ $\in$ ListMessages}
     CO\_Receive(m\'~})
       \ENDFOR
       \ELSE
       \STATE ListMessages.add(m)
       \ENDIF
\end{algorithmic}
\end{algorithm}


\subsection{Question 5 Analize the algorithm}
 blah  blah  blah  blah  blah  blah  blah  blah 

\begin{algorithm}
\caption{\textbf{Q5}(Message m)}
/* \textit{ Au moment où on reçoit un message */}
\label{algo:Q5}
\begin{algorithmic}[1]
\end{algorithmic}
\end{algorithm}



 blah  blah  blah  blah 



\listoffigures
% Pour finir l'interligne de 1,5
%\bibliographystyle{alpha}
%\bibliography{biblio.bib}


\end{document}
